\begin{abstract}
同策略蒸馏（OPD）已逐渐成为将教师模型能力迁移至学生模型的热门方法。在本工作中，我们提出一个关键的研究问题：同策略采样是否总是对任意教师-学生配对的蒸馏有益？我们发现，一种简单的替代方案——半同策略蒸馏（Semi-OPD），即从初始学生生成的离线轨迹中进行蒸馏，通常在准确率和训练效率方面均优于OPD。在涵盖1.5B至235B参数的17组教师-学生配对中，Semi-OPD在14种情况下表现更优，准确率最高提升达+13.6\%，训练速度最快提升11.4$\times$。  
我们进一步发现，OPD与Semi-OPD的选择取决于初始教师与学生之间的对齐程度，该对齐程度可通过输出token重叠率来量化：仅当两者高度对齐且重叠率较高时，OPD才具有优势。  
我们的深入分析表明，有效的蒸馏需要同时针对学生和教师保持同策略性。对于对齐度较差的配对，随着上下文长度的增长，学生采样轨迹会逐渐偏离教师策略，从而削弱蒸馏信号。相比之下，Semi-OPD通常更具稳定性，因为它在较短上下文中进行蒸馏，同时覆盖完整轨迹，并使学生接触到更多教师偏好的token。  
除了提出Semi-OPD作为一种高效替代方案外，我们的工作还促使社区重新思考何时使用OPD，以及如何研究具有实际意义的教师-学生配对的更强版本的OPD。


\end{abstract}

